\documentclass[sigconf]{acmart}

\usepackage{booktabs}
\usepackage{graphicx}
\usepackage{multirow}

% Remove copyright information
\settopmatter{printacmref=false}
\renewcommand\footnotetextcopyrightpermission[1]{}
\pagestyle{plain}

\begin{document}

\title{Comparative Analysis of Supervised Machine Learning Models for Flight Delay Prediction}

\author{Your Name}
\affiliation{%
  \institution{Your University}
  \city{Your City}
  \country{Your Country}
}
\email{your.email@university.edu}

\begin{abstract}
This paper presents a comprehensive comparative analysis of nine supervised machine learning models for flight delay prediction. Using a dataset of 3 million flight records with 20 features, I evaluate K-Nearest Neighbors, Linear Regression, Logistic Regression, Perceptron, Artificial Neural Networks, Decision Trees, Random Forests, Naive Bayes, and Support Vector Machines. The experimental results demonstrate that the Artificial Neural Network achieves the best overall performance with 96.31\% accuracy, 88.70\% F1-Score, and 98.35\% ROC-AUC, outperforming all other models across all evaluation metrics. I provide detailed analysis of each model's strengths and weaknesses, conduct comparative analysis of the top three performers, and justify the selection of Artificial Neural Networks as the optimal model for this classification task.
\end{abstract}

\maketitle

\section{Introduction}

Flight delay prediction is a critical problem in aviation operations, affecting millions of passengers annually and causing significant economic disruptions. Accurate prediction of flight delays enables better resource allocation, improved passenger experience, and optimized airline operations. This study applies and compares nine different supervised machine learning models to predict whether a flight will be delayed based on various operational and temporal features.

The primary objective of this work is to conduct a systematic comparative analysis of these models, evaluating their performance on a large-scale real-world dataset. I address a binary classification task where the target variable indicates whether a flight is delayed (IS\_DELAYED = 1) or not delayed (IS\_DELAYED = 0). Through rigorous experimentation and analysis, I identify the most effective model for this prediction task.

\section{Dataset Description}

\subsection{Dataset Overview}
The dataset consists of 3,000,000 flight records collected from airline operations. After preprocessing, the dataset contains 20 features and 1 binary target variable. The dataset exhibits significant class imbalance with approximately 82.8\% non-delayed flights (2,484,711 samples) and 17.2\% delayed flights (515,289 samples).

\subsection{Features}
The 20 features used in this analysis include:
\begin{itemize}
    \item \textbf{Categorical/Numerical identifiers:} AIRLINE\_CODE, FL\_NUMBER, ORIGIN, DEST
    \item \textbf{Delay metrics:} DEP\_DELAY
    \item \textbf{Operational metrics:} TAXI\_OUT, WHEELS\_OFF, TAXI\_IN, DIVERTED
    \item \textbf{Flight characteristics:} CRS\_ELAPSED\_TIME, DISTANCE
    \item \textbf{Temporal features:} YEAR, MONTH, DAY\_OF\_WEEK
    \item \textbf{Time features:} CRS\_DEP\_TIME\_HOUR, CRS\_DEP\_TIME\_MINUTE, DEP\_TIME\_HOUR, DEP\_TIME\_MINUTE, CRS\_ARR\_TIME\_HOUR, CRS\_ARR\_TIME\_MINUTE
\end{itemize}

\subsection{Data Preprocessing}
The preprocessing pipeline included the following steps:

\textbf{Missing Value Handling:} The dataset contained 470,956 missing values across 9 features. I applied domain-specific imputation strategies:
\begin{itemize}
    \item Features related to cancelled or diverted flights (DEP\_DELAY, TAXI\_OUT, WHEELS\_OFF, TAXI\_IN, DEP\_TIME\_HOUR, DEP\_TIME\_MINUTE) were filled with 0, as these values represent flights where actual operational metrics do not exist.
    \item Scheduled time features (CRS\_ELAPSED\_TIME, CRS\_ARR\_TIME\_HOUR, CRS\_ARR\_TIME\_MINUTE) were filled with median values, as missing values here likely represent data entry errors.
\end{itemize}

\textbf{Feature Selection:} I removed the FL\_DATE column as the temporal information was already captured through the YEAR, MONTH, and DAY\_OF\_WEEK features, avoiding redundancy.

\textbf{Train-Test Split:} I performed an 80-20 stratified split, resulting in 2,400,000 training samples and 600,000 testing samples. Stratification ensures both sets maintain the original class distribution.

\section{Experimental Setup}

\subsection{Evaluation Metrics}
I evaluate all models using the following standard classification metrics:
\begin{itemize}
    \item \textbf{Accuracy:} The proportion of correct predictions among total predictions
    \item \textbf{Precision:} The proportion of true positive predictions among all positive predictions
    \item \textbf{Recall:} The proportion of true positive predictions among all actual positive cases
    \item \textbf{F1-Score:} The harmonic mean of precision and recall, providing a balanced metric
    \item \textbf{ROC-AUC:} Area under the Receiver Operating Characteristic curve, measuring the model's discrimination ability across all classification thresholds
\end{itemize}

Additionally, I record training and prediction times to assess computational efficiency, though this is not a primary selection criterion for this academic study.

\subsection{Implementation Details}
All models were implemented using scikit-learn (version 1.x) in Python 3.11. I used default hyperparameters for all models to ensure fair comparison and avoid optimization bias toward any particular algorithm. Feature scaling was applied when required by specific algorithms (KNN, Perceptron, Logistic Regression, Neural Networks, and SVM).

The experiments were conducted on a standard desktop computer, and all models were trained on the full training set of 2.4 million samples, except for KNN where computational constraints were noted.

\section{Model Descriptions and Results}

\subsection{K-Nearest Neighbors (KNN)}

\subsubsection{Experimental Setup}
I implemented KNN with the following configuration:
\begin{itemize}
    \item \textbf{Parameters:} k=5 neighbors (default), Euclidean distance metric, uniform weights
    \item \textbf{Preprocessing:} StandardScaler for feature normalization
    \item \textbf{Implementation:} KNeighborsClassifier with parallel processing enabled
\end{itemize}

\subsubsection{Performance Evaluation}
\begin{table}[h]
\centering
\caption{KNN Performance Metrics}
\begin{tabular}{lr}
\toprule
Metric & Value \\
\midrule
Accuracy & 0.9250 \\
Precision & 0.9107 \\
Recall & 0.6249 \\
F1-Score & 0.7412 \\
ROC-AUC & 0.9234 \\
Training Time & 0.17s \\
Prediction Time & 702.21s \\
\bottomrule
\end{tabular}
\end{table}

\subsubsection{Analysis}
\textbf{Strengths:}
\begin{itemize}
    \item Simple, non-parametric algorithm requiring minimal training
    \item Makes no assumptions about underlying data distribution
    \item Achieves good precision (91.07\%) when predicting delays
    \item Intuitive decision-making based on similar instances
\end{itemize}

\textbf{Weaknesses:}
\begin{itemize}
    \item Extremely slow prediction time (11.7 minutes for 600K samples), making it impractical for large-scale applications
    \item Low recall (62.49\%), failing to detect 37.5\% of actual delays
    \item Memory-intensive as it stores the entire training dataset
    \item Suffers from the curse of dimensionality with 20 features
    \item Performance degrades with increasing dataset size
\end{itemize}

\textbf{Suitability for the Problem:} KNN demonstrates poor suitability for this flight delay prediction task. While it achieves reasonable accuracy, the computational cost is prohibitive, and the low recall means many actual delays go undetected. The lazy learning approach does not scale well to datasets of this size.

\subsection{Linear Regression}

\subsubsection{Experimental Setup}
I applied Linear Regression with the following setup:
\begin{itemize}
    \item \textbf{Parameters:} Ordinary least squares estimation with default settings
    \item \textbf{Preprocessing:} StandardScaler for feature normalization
    \item \textbf{Classification Adaptation:} Applied threshold at 0.5 (predictions $\geq$ 0.5 classified as delayed)
    \item \textbf{Note:} Linear regression is designed for continuous target prediction but adapted here for binary classification
\end{itemize}

\subsubsection{Performance Evaluation}
\begin{table}[h]
\centering
\caption{Linear Regression Performance Metrics}
\begin{tabular}{lr}
\toprule
Metric & Value \\
\midrule
Accuracy & 0.8780 \\
Precision & 0.9907 \\
Recall & 0.2925 \\
F1-Score & 0.4516 \\
ROC-AUC & 0.9538 \\
Training Time & 0.65s \\
Prediction Time & 0.02s \\
\bottomrule
\end{tabular}
\end{table}

\subsubsection{Analysis}
\textbf{Strengths:}
\begin{itemize}
    \item Very fast training (0.65s) and prediction (0.02s)
    \item Extremely high precision (99.07\%), producing minimal false positives
    \item Good ROC-AUC (95.38\%), indicating strong probabilistic discrimination
    \item Simple, interpretable linear model
    \item Low computational requirements
\end{itemize}

\textbf{Weaknesses:}
\begin{itemize}
    \item Critically low recall (29.25\%), missing 70.75\% of actual delays
    \item Not designed for binary classification tasks
    \item Fixed threshold may not be optimal for imbalanced data
    \item Poorest F1-Score (0.4516) among all evaluated models
    \item Produces unbounded outputs requiring post-processing
\end{itemize}

\textbf{Suitability for the Problem:} Linear Regression demonstrates poor suitability for this classification task. Despite impressive precision and good ROC-AUC, the critically low recall renders it unsuitable for practical delay prediction. Missing 70\% of actual delays would severely limit its utility in real-world applications.

\subsection{Logistic Regression}

\subsubsection{Experimental Setup}
I configured Logistic Regression as follows:
\begin{itemize}
    \item \textbf{Parameters:} Maximum 1000 iterations, L2 regularization (default), lbfgs solver
    \item \textbf{Preprocessing:} StandardScaler for feature normalization
    \item \textbf{Implementation:} LogisticRegression with parallel processing enabled
\end{itemize}

\subsubsection{Performance Evaluation}
\begin{table}[h]
\centering
\caption{Logistic Regression Performance Metrics}
\begin{tabular}{lr}
\toprule
Metric & Value \\
\midrule
Accuracy & 0.9584 \\
Precision & 0.9209 \\
Recall & 0.8288 \\
F1-Score & 0.8724 \\
ROC-AUC & 0.9772 \\
Training Time & 4.31s \\
Prediction Time & 0.02s \\
\bottomrule
\end{tabular}
\end{table}

\subsubsection{Analysis}
\textbf{Strengths:}
\begin{itemize}
    \item Excellent overall performance across all metrics
    \item Well-balanced precision (92.09\%) and recall (82.88\%)
    \item Very fast training and prediction times
    \item Highly interpretable with meaningful feature coefficients
    \item Provides probabilistic outputs suitable for threshold adjustment
    \item Handles class imbalance effectively
    \item Robust and stable performance
\end{itemize}

\textbf{Weaknesses:}
\begin{itemize}
    \item Assumes linear decision boundary, potentially missing non-linear patterns
    \item May underperform when feature interactions are complex
    \item Requires feature scaling for optimal performance
\end{itemize}

\textbf{Suitability for the Problem:} Logistic Regression demonstrates excellent suitability for flight delay prediction. It achieves strong performance across all metrics while maintaining interpretability and computational efficiency. The model provides a solid baseline and represents the best-performing linear classifier in this study.

\subsection{Perceptron Learning}

\subsubsection{Experimental Setup}
I implemented the Perceptron algorithm with:
\begin{itemize}
    \item \textbf{Parameters:} Default parameters with random\_state=42 for reproducibility
    \item \textbf{Preprocessing:} StandardScaler for feature normalization
    \item \textbf{Implementation:} Perceptron with parallel processing enabled
\end{itemize}

\subsubsection{Performance Evaluation}
\begin{table}[h]
\centering
\caption{Perceptron Performance Metrics}
\begin{tabular}{lr}
\toprule
Metric & Value \\
\midrule
Accuracy & 0.9278 \\
Precision & 0.7816 \\
Recall & 0.8049 \\
F1-Score & 0.7930 \\
ROC-AUC & 0.9535 \\
Training Time & 2.09s \\
Prediction Time & 0.01s \\
\bottomrule
\end{tabular}
\end{table}

\subsubsection{Analysis}
\textbf{Strengths:}
\begin{itemize}
    \item Very fast training (2.09s) and prediction (0.01s)
    \item Good recall (80.49\%), detecting most delays
    \item Simple online learning algorithm
    \item Low computational and memory requirements
    \item Good ROC-AUC (95.35\%)
\end{itemize}

\textbf{Weaknesses:}
\begin{itemize}
    \item Lower precision (78.16\%) compared to top models, producing more false positives
    \item Sensitive to feature scaling and training sample order
    \item May not converge on non-linearly separable data
    \item Less stable than Logistic Regression
    \item No probabilistic interpretation without modifications
\end{itemize}

\textbf{Suitability for the Problem:} The Perceptron shows moderate suitability for this task. While computationally efficient with good recall, the lower precision results in more false alarms. It represents a simpler alternative to Logistic Regression but with reduced performance.

\subsection{Artificial Neural Networks (Multi-Layer Perceptron)}

\subsubsection{Experimental Setup}
I designed and trained a neural network with the following architecture:
\begin{itemize}
    \item \textbf{Architecture:} 
    \begin{itemize}
        \item Input layer: 20 neurons (one per feature)
        \item Hidden layer: 100 neurons with ReLU activation
        \item Output layer: 2 neurons with softmax activation
    \end{itemize}
    \item \textbf{Parameters:} Maximum 300 iterations, Adam optimizer (default learning rate 0.001), batch-based training
    \item \textbf{Preprocessing:} StandardScaler for feature normalization (critical for neural networks)
    \item \textbf{Implementation:} MLPClassifier from scikit-learn
    \item \textbf{Convergence:} The model converged in 27 iterations
\end{itemize}

\subsubsection{Performance Evaluation}
\begin{table}[h]
\centering
\caption{Neural Network Performance Metrics}
\begin{tabular}{lr}
\toprule
Metric & Value \\
\midrule
Accuracy & \textbf{0.9631} \\
Precision & \textbf{0.9349} \\
Recall & \textbf{0.8438} \\
F1-Score & \textbf{0.8870} \\
ROC-AUC & \textbf{0.9835} \\
Training Time & 141.68s \\
Prediction Time & 0.76s \\
Iterations & 27 \\
\bottomrule
\end{tabular}
\end{table}

\subsubsection{Analysis}
\textbf{Strengths:}
\begin{itemize}
    \item \textbf{Best overall performance across all evaluation metrics}
    \item Highest accuracy (96.31\%), surpassing all other models
    \item Highest F1-Score (88.70\%), indicating superior balanced performance
    \item Highest recall (84.38\%), detecting more delays than any other model
    \item Highest ROC-AUC (98.35\%), demonstrating excellent discrimination ability
    \item Highest precision (93.49\%), minimizing false positives
    \item Capable of learning complex non-linear patterns in the data
    \item Efficient convergence in only 27 iterations
    \item Flexible architecture that can be expanded for more complex patterns
\end{itemize}

\textbf{Weaknesses:}
\begin{itemize}
    \item Longer training time (141.68s) compared to linear models
    \item Black-box nature makes interpretation more challenging
    \item Requires careful feature scaling
    \item Potential for overfitting without proper regularization
    \item Sensitive to initialization and hyperparameter choices
\end{itemize}

\textbf{Suitability for the Problem:} The Artificial Neural Network demonstrates excellent suitability for flight delay prediction. It achieves the best performance across all metrics, indicating it has successfully learned complex patterns in the flight data that linear models cannot capture. The training time of approximately 2.5 minutes is acceptable for a model that will be trained once and used for many predictions. The superior recall means it detects more actual delays, which is crucial for operational planning.

\subsection{Decision Trees}

\subsubsection{Experimental Setup}
I implemented Decision Trees with:
\begin{itemize}
    \item \textbf{Parameters:} Default parameters (no maximum depth limit), random\_state=42
    \item \textbf{Preprocessing:} No scaling required (tree-based methods are scale-invariant)
    \item \textbf{Implementation:} DecisionTreeClassifier from scikit-learn
\end{itemize}

\subsubsection{Performance Evaluation}
\begin{table}[h]
\centering
\caption{Decision Tree Performance Metrics}
\begin{tabular}{lr}
\toprule
Metric & Value \\
\midrule
Accuracy & 0.9336 \\
Precision & 0.7984 \\
Recall & 0.8204 \\
F1-Score & 0.8093 \\
ROC-AUC & 0.8887 \\
Training Time & 27.85s \\
Prediction Time & 0.22s \\
Tree Depth & 49 \\
Number of Leaves & 96,898 \\
\bottomrule
\end{tabular}
\end{table}

\subsubsection{Analysis}
\textbf{Strengths:}
\begin{itemize}
    \item No preprocessing or feature scaling required
    \item Naturally interpretable structure (when not too deep)
    \item Handles non-linear relationships effectively
    \item Good recall (82.04\%), detecting most delays
    \item Can identify important features through the tree structure
\end{itemize}

\textbf{Weaknesses:}
\begin{itemize}
    \item Lowest ROC-AUC (88.87\%) among classification models, indicating weaker probabilistic predictions
    \item Extremely deep tree (49 levels) with 96,898 leaves indicates severe overfitting
    \item Poor generalization due to overfitting on training data
    \item Lower precision (79.84\%) compared to top models
    \item Unstable - small changes in data can produce very different trees
    \item Not practically interpretable due to excessive depth
\end{itemize}

\textbf{Suitability for the Problem:} Decision Trees show poor to moderate suitability for this task. The excessive tree depth clearly indicates overfitting, where the model has memorized training data rather than learning generalizable patterns. The significantly lower ROC-AUC compared to other models confirms poor generalization. Standalone Decision Trees are not recommended for this problem.

\subsection{Random Forests}

\subsubsection{Experimental Setup}
I configured Random Forests as follows:
\begin{itemize}
    \item \textbf{Parameters:} 100 trees (n\_estimators=100), default maximum depth, random\_state=42
    \item \textbf{Preprocessing:} No scaling required
    \item \textbf{Implementation:} RandomForestClassifier with parallel processing enabled
\end{itemize}

\subsubsection{Performance Evaluation}
\begin{table}[h]
\centering
\caption{Random Forest Performance Metrics}
\begin{tabular}{lr}
\toprule
Metric & Value \\
\midrule
Accuracy & 0.9575 \\
Precision & 0.9203 \\
Recall & 0.8237 \\
F1-Score & 0.8693 \\
ROC-AUC & 0.9781 \\
Training Time & 50.64s \\
Prediction Time & 1.94s \\
\bottomrule
\end{tabular}
\end{table}

\subsubsection{Analysis}
\textbf{Strengths:}
\begin{itemize}
    \item Excellent overall performance, third-best F1-Score (86.93\%)
    \item Dramatic improvement over single Decision Tree through ensemble averaging
    \item High ROC-AUC (97.81\%), excellent discrimination ability
    \item No preprocessing required
    \item Robust to overfitting compared to single trees
    \item Provides feature importance rankings
    \item Handles non-linear relationships naturally
\end{itemize}

\textbf{Weaknesses:}
\begin{itemize}
    \item Moderate training time (50.64s)
    \item Slower prediction time (1.94s) compared to linear models
    \item Less interpretable than single decision trees
    \item Memory-intensive, storing 100 complete trees
    \item Slightly lower recall (82.37\%) than Neural Network
\end{itemize}

\textbf{Suitability for the Problem:} Random Forests demonstrate very good suitability for flight delay prediction. The ensemble approach successfully addresses the overfitting issues of single Decision Trees while maintaining strong performance. It ranks among the top three models and would serve as a reliable alternative to neural networks.

\subsection{Naive Bayes}

\subsubsection{Experimental Setup}
I implemented Gaussian Naive Bayes with:
\begin{itemize}
    \item \textbf{Variant:} Gaussian Naive Bayes (assumes features follow normal distributions)
    \item \textbf{Parameters:} Default parameters
    \item \textbf{Preprocessing:} No scaling (works with original feature distributions)
    \item \textbf{Implementation:} GaussianNB from scikit-learn
    \item \textbf{Note:} Initial experiments with scaled features produced very poor results; the model performs best with original distributions
\end{itemize}

\subsubsection{Performance Evaluation}
\begin{table}[h]
\centering
\caption{Naive Bayes Performance Metrics}
\begin{tabular}{lr}
\toprule
Metric & Value \\
\midrule
Accuracy & 0.9370 \\
Precision & 0.8225 \\
Recall & 0.8073 \\
F1-Score & 0.8148 \\
ROC-AUC & 0.9532 \\
Training Time & 0.43s \\
Prediction Time & 0.20s \\
\bottomrule
\end{tabular}
\end{table}

\subsubsection{Analysis}
\textbf{Strengths:}
\begin{itemize}
    \item Extremely fast training (0.43s) and prediction (0.20s)
    \item Simple, probabilistic model with clear assumptions
    \item Good ROC-AUC (95.32\%), indicating strong discrimination
    \item Minimal computational and memory requirements
    \item Works well with original feature distributions
\end{itemize}

\textbf{Weaknesses:}
\begin{itemize}
    \item Strong independence assumption is violated (flight features are correlated)
    \item Lower precision (82.25\%) and recall (80.73\%) compared to top models
    \item Performance significantly degraded when features were scaled
    \item Unrealistic assumption that all features are independent
    \item Cannot capture feature interactions
\end{itemize}

\textbf{Suitability for the Problem:} Naive Bayes shows moderate suitability for this task. While computationally efficient, the independence assumption is clearly inappropriate for flight data where features like departure delay, taxi times, and distance are inherently correlated. The model achieves reasonable but not competitive performance.

\subsection{Support Vector Machines (Linear SVM)}

\subsubsection{Experimental Setup}
I implemented Linear SVM with:
\begin{itemize}
    \item \textbf{Kernel:} Linear kernel
    \item \textbf{Parameters:} Maximum 1000 iterations, random\_state=42
    \item \textbf{Preprocessing:} StandardScaler for normalization, followed by CalibratedClassifierCV for probability estimates
    \item \textbf{Implementation:} LinearSVC from scikit-learn
    \item \textbf{Important Note:} I initially attempted standard SVM with RBF kernel but terminated training after 75 minutes without convergence. The quadratic computational complexity of RBF kernel SVM makes it infeasible for datasets of this scale (2.4 million samples). Linear SVM was chosen as a practical alternative.
\end{itemize}

\subsubsection{Performance Evaluation}
\begin{table}[h]
\centering
\caption{Linear SVM Performance Metrics}
\begin{tabular}{lr}
\toprule
Metric & Value \\
\midrule
Accuracy & 0.9582 \\
Precision & 0.9206 \\
Recall & 0.8280 \\
F1-Score & 0.8719 \\
ROC-AUC & 0.9770 \\
Training Time & 28.89s \\
Prediction Time & 0.11s \\
\bottomrule
\end{tabular}
\end{table}

\subsubsection{Analysis}
\textbf{Strengths:}
\begin{itemize}
    \item Excellent performance, virtually identical to Logistic Regression
    \item Robust to outliers through margin maximization
    \item Fast prediction time (0.11s)
    \item Effective on high-dimensional data
    \item Strong theoretical foundation based on statistical learning theory
    \item Good generalization through margin maximization
\end{itemize}

\textbf{Weaknesses:}
\begin{itemize}
    \item Limited to linear decision boundaries (non-linear kernels computationally prohibitive)
    \item RBF kernel version infeasible for this dataset size
    \item Requires feature scaling
    \item Calibration step needed for probability estimates
    \item Less interpretable than Logistic Regression
    \item Longer training time (28.89s) than Logistic Regression (4.31s) without performance benefit
\end{itemize}

\textbf{Suitability for the Problem:} Linear SVM demonstrates very good suitability and achieves performance nearly identical to Logistic Regression. However, it offers no significant advantages while being computationally more expensive and less interpretable. The inability to use non-linear kernels at this scale is a notable limitation.

\section{Comparative Analysis of Top 3 Models}

Based on F1-Score, the three best-performing models are:
\begin{enumerate}
    \item Artificial Neural Network (MLP): F1-Score = 0.8870
    \item Logistic Regression: F1-Score = 0.8724
    \item Linear SVM: F1-Score = 0.8719
\end{enumerate}

\subsection{Comprehensive Performance Comparison}

\begin{table*}[t]
\centering
\caption{Detailed Comparison of Top 3 Models}
\begin{tabular}{lrrrrrrr}
\toprule
Model & Accuracy & Precision & Recall & F1-Score & ROC-AUC & Train Time & Pred Time \\
\midrule
Neural Network & \textbf{0.9631} & \textbf{0.9349} & \textbf{0.8438} & \textbf{0.8870} & \textbf{0.9835} & 141.68s & 0.76s \\
Logistic Regression & 0.9584 & 0.9209 & 0.8288 & 0.8724 & 0.9772 & \textbf{4.31s} & \textbf{0.02s} \\
Linear SVM & 0.9582 & 0.9206 & 0.8280 & 0.8719 & 0.9770 & 28.89s & 0.11s \\
\bottomrule
\end{tabular}
\end{table*}

\subsection{Detailed Analysis of Performance Differences}

\subsubsection{Neural Network vs. Logistic Regression}

The Neural Network consistently outperforms Logistic Regression across all metrics:

\textbf{Accuracy:} The Neural Network achieves 96.31\% accuracy compared to Logistic Regression's 95.84\%, an improvement of 0.47 percentage points. On the test set of 600,000 samples, this translates to approximately 2,820 additional correct predictions.

\textbf{Precision:} The Neural Network achieves 93.49\% precision versus 92.09\% for Logistic Regression, a gain of 1.40 percentage points. This means the Neural Network produces fewer false positive predictions - specifically, 1,280 fewer false alarms in the test set.

\textbf{Recall:} The most significant improvement is in recall, where the Neural Network achieves 84.38\% compared to 82.88\% for Logistic Regression, an improvement of 1.50 percentage points. This translates to detecting approximately 1,545 additional actual delays that Logistic Regression would miss. In a practical setting, this improved recall is particularly valuable for operational planning.

\textbf{F1-Score:} The Neural Network achieves an F1-Score of 88.70\% versus 87.24\% for Logistic Regression, a difference of 1.46 percentage points. This indicates better balanced performance across both precision and recall.

\textbf{ROC-AUC:} The Neural Network achieves 98.35\% ROC-AUC compared to 97.72\% for Logistic Regression, a gain of 0.63 percentage points. This superior ROC-AUC indicates better probabilistic predictions across all possible classification thresholds, which is valuable for threshold optimization.

\subsubsection{Confusion Matrix Analysis}

\begin{table}[h]
\centering
\caption{Detailed Error Analysis - Top 3 Models}
\begin{tabular}{lrrrr}
\toprule
Model & True Neg & False Pos & False Neg & True Pos \\
\midrule
Neural Network & 490,889 & 6,053 & 16,098 & 86,960 \\
Logistic Regression & 489,609 & 7,333 & 17,643 & 85,415 \\
Linear SVM & 489,578 & 7,364 & 17,722 & 85,336 \\
\bottomrule
\end{tabular}
\end{table}

The Neural Network demonstrates superior performance in both error types:
\begin{itemize}
    \item \textbf{False Positives:} The Neural Network produces 6,053 false positives, 1,280 fewer than Logistic Regression (17.5\% reduction) and 1,311 fewer than Linear SVM (17.8\% reduction).
    \item \textbf{False Negatives:} The Neural Network produces 16,098 false negatives, 1,545 fewer than Logistic Regression (8.8\% reduction) and 1,624 fewer than Linear SVM (9.2\% reduction).
    \item \textbf{True Positives:} The Neural Network correctly identifies 86,960 delays, 1,545 more than Logistic Regression and 1,624 more than Linear SVM.
\end{itemize}

This balanced improvement across both error types is significant. Fewer false negatives mean more actual delays are detected, while fewer false positives mean less unnecessary alert fatigue.

\subsubsection{Logistic Regression vs. Linear SVM}

These two models produce remarkably similar results:
\begin{itemize}
    \item \textbf{F1-Score difference:} Only 0.05 percentage points (0.8724 vs. 0.8719)
    \item \textbf{Accuracy difference:} Only 0.02 percentage points (95.84\% vs. 95.82\%)
    \item \textbf{All metrics differ by less than 0.1 percentage points}
\end{itemize}

Despite their algorithmic differences, both models converge to nearly identical solutions. This makes sense theoretically: both are linear classifiers that learn decision hyperplanes, with Logistic Regression minimizing log loss and Linear SVM maximizing margin. On large, well-separated datasets, these objectives lead to similar solutions.

The primary practical difference is computational efficiency: Logistic Regression trains 6.7 times faster (4.31s vs. 28.89s) and predicts 5.5 times faster (0.02s vs. 0.11s) while achieving virtually identical performance.

\subsection{Analysis of Variability and Underlying Reasons}

\subsubsection{Why Does the Neural Network Outperform Linear Models?}

The consistent superiority of the Neural Network suggests the presence of non-linear patterns in the flight delay data that linear models cannot capture. Several factors explain this:

\textbf{1. Non-linear Feature Interactions:} The hidden layer with 100 neurons and ReLU activation functions enables the network to learn complex interactions between features. For example, the relationship between departure delay and flight delay might differ based on time of day, airline, or distance - interactions that linear models cannot represent.

\textbf{2. Adaptive Feature Representation:} The neural network effectively learns a new feature representation in its hidden layer. The 100 hidden neurons create 100 new features that are non-linear combinations of the original 20 features, allowing the model to discover patterns that are not linearly separable in the original feature space.

\textbf{3. Flexibility in Decision Boundaries:} While linear models are constrained to linear decision boundaries (hyperplanes), the neural network can learn arbitrarily complex decision boundaries. This flexibility allows it to better separate delayed and non-delayed flights in regions where the two classes overlap.

\textbf{4. Class Imbalance Handling:} With 82.8\% non-delayed flights, the dataset is imbalanced. The Neural Network's iterative training with mini-batches and adaptive learning may handle this imbalance more effectively than the closed-form or margin-based solutions of linear models.

\subsubsection{Why Is the Performance Gap Relatively Small?}

Despite the Neural Network's architectural advantages, the performance improvement over linear models is modest (1.46 percentage points in F1-Score). This observation is informative:

\textbf{1. Data Is Largely Linearly Separable:} The high performance of linear models (Logistic Regression achieving 87.24\% F1-Score) indicates that much of the prediction task can be solved with linear combinations of features. This suggests that features like departure delay, taxi times, and scheduled times have strong linear relationships with the outcome.

\textbf{2. Limited Non-linear Patterns:} While non-linear patterns exist (evidenced by the Neural Network's superior performance), they represent a relatively small portion of the predictable variance. The majority of delay patterns can be captured through linear relationships.

\textbf{3. Feature Quality:} The preprocessed features already capture much of the relevant information. The 20 features were well-chosen and cleaned, reducing the need for complex non-linear transformations.

\subsubsection{Statistical Significance}

With 600,000 test samples, even small performance differences are statistically significant. The Neural Network's 1.46 percentage point improvement in F1-Score corresponds to correctly classifying approximately 8,760 additional flights. Over millions of predictions in operational settings, this improvement becomes substantial.

\subsection{Computational Considerations}

\textbf{Training Time:} The Neural Network requires 141.68 seconds (approximately 2.4 minutes) to train, compared to 4.31 seconds for Logistic Regression - a 33-fold difference. However, in the context of this academic study and practical applications where models are trained infrequently (perhaps weekly or monthly on new data), this training time is entirely acceptable.

\textbf{Prediction Time:} The Neural Network requires 0.76 seconds to make predictions on 600,000 samples (1.27 microseconds per prediction), compared to 0.02 seconds for Logistic Regression (33 nanoseconds per prediction). For batch prediction scenarios, both times are negligible. Even for real-time predictions, 1.27 microseconds per flight is well within acceptable latency requirements.

\textbf{Model Size and Deployment:} The Neural Network with 100 hidden neurons has approximately 2,200 parameters ((20 × 100) + (100 × 2) + biases), compared to 21 parameters for Logistic Regression. While larger, this model size is still very small by modern standards and poses no deployment challenges.

\section{Model Selection and Justification}

\subsection{Final Recommendation: Artificial Neural Network}

Based on the comprehensive experimental analysis and comparative evaluation, I recommend the \textbf{Artificial Neural Network (Multi-Layer Perceptron)} as the optimal model for flight delay prediction. This recommendation is justified by the following factors:

\subsubsection{1. Superior Predictive Performance}

The Neural Network achieves the best performance across all evaluation metrics:
\begin{itemize}
    \item \textbf{Highest Accuracy (96.31\%):} Correctly classifies more flights than any other model
    \item \textbf{Highest Precision (93.49\%):} Minimizes false alarms
    \item \textbf{Highest Recall (84.38\%):} Detects more actual delays, crucial for operational planning
    \item \textbf{Highest F1-Score (88.70\%):} Best balanced performance
    \item \textbf{Highest ROC-AUC (98.35\%):} Superior discrimination across all thresholds
\end{itemize}

The consistent superiority across all metrics, not just one, demonstrates robust and reliable performance.

\subsubsection{2. Practical Significance of Performance Gains}

While the improvement over Logistic Regression may appear modest (1.46 percentage points in F1-Score), the practical impact is significant:

\begin{itemize}
    \item \textbf{8,760 additional correct predictions} on a 600,000-sample test set
    \item \textbf{1,545 fewer missed delays} (8.8\% reduction in false negatives)
    \item \textbf{1,280 fewer false alarms} (17.5\% reduction in false positives)
\end{itemize}

In an operational context processing millions of flight predictions, these improvements translate to tangible value: better resource allocation, reduced passenger impact, and improved decision-making.

\subsubsection{3. Acceptable Computational Requirements}

The training time of 141.68 seconds (approximately 2.4 minutes) is entirely acceptable for several reasons:

\textbf{One-Time Training:} Machine learning models are typically trained once and then used for numerous predictions. In production settings, a model might be retrained weekly or monthly with updated data. A 2.4-minute training time for a model that will serve millions of predictions is negligible.

\textbf{Academic Context:} In a research or academic setting, training time is not a critical constraint. The priority is achieving the best possible understanding of the data and the most accurate predictions.

\textbf{Prediction Efficiency:} The prediction time of 0.76 seconds for 600,000 samples (1.27 microseconds per prediction) is extremely fast and suitable for both batch and real-time scenarios.

\subsubsection{4. Capacity for Non-linear Pattern Learning}

The Neural Network's architecture enables it to discover and leverage non-linear relationships in the data that linear models cannot represent. The hidden layer effectively learns a transformed feature space where complex patterns become apparent. This capability is evidenced by its superior performance and suggests that flight delay patterns involve interactions and non-linearities beyond simple linear combinations.

\subsubsection{5. Scalability and Extensibility}

The Neural Network architecture used here (20-100-2) is relatively simple and can be easily extended:
\begin{itemize}
    \item Additional hidden layers could be added for more complex pattern learning
    \item Hidden layer size could be increased
    \item Advanced techniques like dropout, batch normalization, or different activation functions could be incorporated
    \item The architecture provides a foundation for further experimentation and improvement
\end{itemize}

\subsubsection{6. Robust Generalization}

The Neural Network achieved excellent test set performance while converging in only 27 iterations, well below the maximum of 300. This suggests:
\begin{itemize}
    \item The model is not overfitting despite its increased complexity
    \item The optimization process found a good solution efficiently
    \item The model is likely to generalize well to new, unseen flight data
\end{itemize}

\subsection{Trade-offs and Considerations}

I acknowledge the following trade-offs in selecting the Neural Network:

\textbf{Interpretability:} Unlike Logistic Regression, which provides interpretable feature coefficients, the Neural Network operates as a black-box. Understanding exactly why it makes specific predictions is more challenging. However, techniques like feature importance analysis, SHAP values, or attention mechanisms could be applied if interpretability becomes critical.

\textbf{Training Time:} The 33× longer training time compared to Logistic Regression is a trade-off for superior performance. Given that training occurs infrequently, this is an acceptable trade-off in both academic and practical contexts.

\textbf{Hyperparameter Sensitivity:} Neural Networks can be sensitive to hyperparameter choices (learning rate, architecture, etc.). However, using default parameters, the model achieved excellent results, suggesting robustness in this application.

\subsection{When Alternative Models Might Be Preferred}

While I recommend the Neural Network as the primary choice, alternative scenarios where other models might be preferred include:

\textbf{Logistic Regression might be preferred if:}
\begin{itemize}
    \item Strict interpretability requirements exist (e.g., regulatory compliance)
    \item Minimal computational infrastructure is available
    \item A simple baseline is needed for comparison
    \item The ~1.5\% performance difference is acceptable
\end{itemize}

\textbf{Random Forest might be preferred if:}
\begin{itemize}
    \item Feature importance rankings are needed
    \item No feature scaling or preprocessing is desired
    \item An ensemble approach is valued for its robustness
\end{itemize}

However, for the primary objective of maximizing predictive accuracy in flight delay prediction, the Artificial Neural Network is the clear choice.

\section{Conclusion}

This study conducted a comprehensive comparative analysis of nine supervised machine learning models for flight delay prediction using a large-scale dataset of 3 million flight records. Through systematic experimentation and evaluation, several important conclusions emerge:

\subsection{Key Findings}

\textbf{1. Performance Hierarchy:} The models rank as follows by F1-Score:
\begin{enumerate}
    \item Artificial Neural Network (88.70\%)
    \item Logistic Regression (87.24\%)
    \item Linear SVM (87.19\%)
    \item Random Forest (86.93\%)
    \item Naive Bayes (81.48\%)
    \item Decision Tree (80.93\%)
    \item Perceptron (79.30\%)
    \item K-Nearest Neighbors (74.12\%)
    \item Linear Regression (45.16\%)
\end{enumerate}

\textbf{2. Neural Network Superiority:} The Artificial Neural Network outperformed all other models across all evaluation metrics, demonstrating its ability to learn complex, non-linear patterns in flight delay data.

\textbf{3. Linear Model Competitiveness:} Despite their simplicity, linear models (Logistic Regression and Linear SVM) achieved strong performance (>87\% F1-Score), indicating that much of the flight delay prediction task is linearly separable.

\textbf{4. Ensemble Benefits:} Random Forests dramatically improved upon single Decision Trees (86.93\% vs. 80.93\% F1-Score), demonstrating the value of ensemble methods in reducing overfitting.

\textbf{5. Computational Feasibility:} Most models trained in reasonable time (<3 minutes), with the exception of K-Nearest Neighbors' prediction time and standard SVM's training infeasibility at this scale.

\subsection{Final Recommendation}

I recommend the \textbf{Artificial Neural Network} as the optimal model for flight delay prediction. With 96.31\% accuracy, 88.70\% F1-Score, and 98.35\% ROC-AUC, it provides the best predictive performance. The training time of 2.4 minutes is entirely acceptable for a model that will be trained infrequently and used for millions of predictions. The superior recall (84.38\%) is particularly valuable, as detecting actual delays is critical for operational planning and passenger communication.

\subsection{Future Work}

Several directions for future research could further improve flight delay prediction:

\textbf{1. Hyperparameter Optimization:} This study used default parameters for fair comparison. Systematic hyperparameter tuning (e.g., grid search, random search, or Bayesian optimization) could improve all models, particularly the Neural Network.

\textbf{2. Advanced Neural Architectures:} Deeper networks, different activation functions (e.g., LeakyReLU, ELU), regularization techniques (dropout, L1/L2), and normalization layers (batch normalization) could enhance performance.

\textbf{3. Feature Engineering:} Creating additional features such as:
\begin{itemize}
    \item Interaction terms between key features
    \item Temporal features (rush hour indicators, seasonal patterns)
    \item Airport-specific congestion metrics
    \item Airline-specific delay history
    \item Weather data integration
\end{itemize}

\textbf{4. Addressing Class Imbalance:} Techniques such as SMOTE (Synthetic Minority Over-sampling Technique), class weights, or threshold adjustment could better handle the 82.8\% / 17.2\% class imbalance.

\textbf{5. Ensemble Methods:} Combining multiple top models (Neural Network, Logistic Regression, Random Forest) through stacking or voting could potentially achieve even better performance than any single model.

\textbf{6. Deep Learning Approaches:} More sophisticated deep learning architectures (e.g., LSTM for temporal patterns, attention mechanisms) could be explored if additional temporal sequence data becomes available.

\textbf{7. External Data Integration:} Incorporating real-time weather data, air traffic control information, and airport congestion metrics could significantly enhance prediction accuracy.

\textbf{8. Interpretability Analysis:} Applying explainable AI techniques (SHAP values, LIME, attention visualization) to understand what the Neural Network learns could provide valuable insights into delay causation.

\subsection{Contribution}

This work provides a thorough, systematic comparison of supervised learning approaches for flight delay prediction on a large-scale real-world dataset. The findings demonstrate that while linear models provide strong baseline performance, neural networks can achieve superior results by learning non-linear patterns in the data. The comprehensive analysis, including detailed performance metrics, error analysis, and practical considerations, provides a solid foundation for both academic understanding and potential practical deployment of flight delay prediction systems.

\end{document}
